% ======================================================================
% 予想4.1 の定理化 — 卒論用の清書 v3（2巡目査読 P9v2 反映版）
% v2 からの変更点（2巡目査読 fable_points2.md §4 の 1〜8 に対応．9 は proof_notes.md の更新）:
%  (1) 命題環境が §2 全体を包含していた構造破損を修復（命題を §1 内で閉じる）
%  (2) C(α) の定義を 4/(∫√(F^{-1})')^2 に修正（分子 α は誤り）．命題1.3 を定数付きの厳密な形に
%  (3) 共同極限を仮定 (H)（Zador 漸近の α 一様性）を明示した系 1.4 に降格し，厳密冪裾で (H) が成立する例を追加
%  (4) 注意（β=1）の regime 混在を解消（固定 α の式と α↓0 の付記を分離）
%  (5) 「歪み最適配置 = F^{-1}(αU) の L^1 最適量子化」の同一視を補題として証明（補題3.2）．
%      Zador の条件が定義1.1 から従うことを補題化（補題3.3）し，主定理の外部仮定から削除．Graf–Luschgy Thm 6.2 を明記
%  (6) 補題3.1 の添字・部分区間・下界の (1-ε) 因子を修正．隠れ定数の α 非依存性を明記
%  (7) 定義1.1 に (F^{-1})'(α)>0 を追加，補題2.3 に C^1 を明記，非増加性が α≤1/2 を課すことを注記
%  (8) Karamata・Potter の文献（Bingham–Goldie–Teugels）を明記．編集履歴コメントを削除
% 表記は docs/glossary.md に従う（区間／下側裾／原子／t分布 等）
% コンパイル: platex + dvipdfmx（jsarticle）
% ======================================================================
\documentclass[a4paper,10pt]{jsarticle}
\usepackage{amsmath,amssymb,amsthm}
\usepackage[margin=25mm]{geometry}

\theoremstyle{definition}
\newtheorem{definition}{定義}[section]
\newtheorem{proposition}{命題}[section]
\newtheorem{lemma}{補題}[section]
\newtheorem{theorem}{定理}[section]
\newtheorem{corollary}{系}[section]
\newtheorem{conjecture}{予想}[section]
\newtheorem{remark}{注意}[section]
\newtheorem{example}{例}[section]
\makeatletter
\renewenvironment{proof}[1][\proofname]{\par
  \pushQED{\qed}%
  \normalfont \topsep6\p@\@plus6\p@\relax
  \trivlist
  \item[\hskip\labelsep\textbf{#1}]\ignorespaces
}{%
  \popQED\endtrivlist\@endpefalse
}
\renewcommand{\proofname}{証明}
\makeatother

\newcommand{\CVaR}{\mathrm{CVaR}}
\newcommand{\Wone}{W_1}
\expandafter\let\csname Finv\endcsname\undefined
\newcommand{\Finv}{F^{-1}}
\newcommand{\Fhinv}{\hat{F}^{-1}}
\newcommand{\R}{\mathbb{R}}
\newcommand{\E}{\mathbb{E}}
\newcommand{\Prob}{\mathbb{P}}
\newcommand{\eN}{e_{N,1}}
\newcommand{\Eunif}{E_{\mathrm{unif}}}
\newcommand{\Eopt}{E_{\mathrm{opt}}}
\newcommand{\Iunif}{I_\alpha^{\mathrm{unif}}}
\newcommand{\Iopt}{I_\alpha^{\mathrm{opt}}}

\title{CVaR 上界汎関数の配置間比較\\
  一様配置と歪み最適配置の漸近スケーリング\\
  （予想4.1 の定理化・査読反映版 v3）}
\author{}
\date{}
\begin{document}
\maketitle

\section{設定と主張}
\label{sec:setting}

本稿では，中間レポート（第4章）で予想として提示した，歪み最適配置による $\CVaR$ 推定誤差上界の改善スケーリングを，正則変動の仮定のもとで定理として証明する．
比較の対象は実際の $\CVaR$ 推定誤差そのものではなく，式\eqref{eq:bound-cvar2}の右辺に現れる\textbf{同一の上界汎関数} $\frac1\alpha I_\alpha$ を，
一様配置と歪み最適配置のそれぞれで評価した値の比である（注意\ref{rem:scope}）．この区別は定理の解釈上重要である．

\paragraph{記法}
$Z$ を損益（大きいほど良い）とし，その分位点関数を $\Finv\colon(0,1)\to\R$（累積分布関数 $F$ の一般化逆関数．非減少・左連続）とする．
条件付きバリューアットリスク（$\CVaR$）の歪み関数 $g_\alpha(\tau)=\min(\tau/\alpha,1)$（$g_\alpha'(\tau)=\alpha^{-1}\mathbb{1}_{[0,\alpha)}(\tau)$）に対し，
誤差上界（中間レポート 命題3.1）は
\begin{equation}
  \bigl|\CVaR_\alpha(Z)-\CVaR_\alpha(\hat Z_N)\bigr|
  \;\le\;\frac1\alpha\int_0^\alpha\bigl|\Finv(\tau)-\Fhinv(\tau)\bigr|\,d\tau
  \;=:\;\frac1\alpha\,I_\alpha(\hat Z_N)
  \label{eq:bound-cvar2}
\end{equation}
である．ここで $\hat Z_N$ は $N$ 個の原子（アトム，atom．確率質量を置く離散的な点）をもつ近似分布，$\Fhinv$ はその分位点関数（階段関数）である．

\textbf{一様配置}（FQF 相当）とは，quantile fraction（累積確率軸 $[0,1]$ の区切り位置）を $\tau_i=i/N$（$i=0,\dots,N$）とし，
区間 $J_i=[\tau_i,\tau_{i+1})$ 上で $\Fhinv(\tau)=\Finv\bigl((i+\tfrac12)/N\bigr)$（各区間の中点における真の分位点値）とおく配置をいう．その上界の評価値を
$\Eunif:=\frac1\alpha\Iunif$，$\Iunif:=I_\alpha(\hat Z_N^{\mathrm{unif}})$ と書く．

\textbf{歪み最適配置}（提案法の理想化．正則化パラメータ $\lambda=0$）とは，$N$ 原子分布の全体にわたる $I_\alpha$ の最小値を達成する配置をいい，
\begin{equation}
  \Iopt:=\min_{\hat Z_N}I_\alpha(\hat Z_N),\qquad \Eopt:=\frac1\alpha\Iopt
  \label{eq:def-opt}
\end{equation}
とおく（最小値の存在と，それが $L^1$ 最適量子化に一致することは補題\ref{lem:identify} で示す）．

\begin{definition}[分位点関数の正則変動]
\label{def:regvar}
$\alpha\in(0,1)$ を固定する．$\Finv\in C^1(0,\alpha]$ かつ $\Finv\in L^1[0,\alpha]$（$\CVaR_\alpha$ が有限値をもつ）とする．
さらに導関数が原点近傍（$\tau\downarrow0$）で正則変動する，すなわち
\begin{equation}
  (\Finv)'(\tau)\sim\tau^{-\beta}\ell(\tau)\qquad(\tau\downarrow0)
  \label{eq:regvar}
\end{equation}
とする．ここで $\beta\in[1,2)$ は発散指数，$\ell$ は緩変動関数（正値・可測で，任意の $c>0$ に対し $\ell(c\tau)/\ell(\tau)\to1$）である．
この漸近関係は原点近傍でのみ課し，$(0,\alpha]$ 全体での厳密な等式は仮定しない．
あわせて $(\Finv)'$ は $(0,\alpha]$ 上で非増加（すなわち $\Finv$ は $(0,\alpha]$ で凹）かつ $(\Finv)'(\alpha)>0$ とする．
非増加性より $(\Finv)'>0$ が $(0,\alpha]$ 全体で成り立ち，$\Finv$ は $(0,\alpha]$ で狭義単調増大である．
\end{definition}

\begin{remark}[仮定の射程]
\label{rem:assumption}
(i) $(\Finv)'$ の非増加性は，対称単峰分布では実質的に $\alpha\le1/2$ を要求する（$\tau>1/2$ で $(\Finv)'$ は増加に転じる）．以下で $\alpha\in(0,1)$ と書くときも，定義\ref{def:regvar} が成立する範囲に限る．
(ii) 式\eqref{eq:regvar}の $\sim$ の収束速度は $\Finv$ に固有の量であり，以下の $\asymp$ 評価に現れる定数はこの収束速度（すなわち $\Finv$ 自体）に依存する．
\end{remark}

\begin{example}
自由度 $\nu>1$ の $t$分布では $(\Finv)'(\tau)\sim\frac{c}{\nu}\tau^{-(1+1/\nu)}$ であり $\beta=1+1/\nu\in(1,2)$（$\nu=3$ で $\beta=4/3$），$\alpha\le1/2$ で非増加性も満たす．
ガウス裾では $\beta=1$ で $(\Finv)'(\tau)\sim\tau^{-1}\bigl(2\ln(1/\tau)\bigr)^{-1/2}$（対数的緩変動）となる．
指数分布や通常の対数正規分布は下側分位点が有限な下端点をもち（$-\infty$ に発散しない），本枠組みの対象外である（注意\ref{rem:beta1} 参照）．
\end{example}

\subsection{主定理（固定 $\alpha$ での $N$ 漸近）}

\begin{theorem}[一様配置と歪み最適配置における上界汎関数の比（固定 $\alpha$）]
\label{thm:alpha}
$\alpha$ を固定し，定義\ref{def:regvar} の仮定が $\beta\in(1,2)$ で成り立つとする．このとき $N\to\infty$ で
\begin{equation}
  \frac{\Eunif}{\Eopt}
  \;\asymp\;C(\alpha)\,N^{\beta-1}\,\ell(1/N),
  \qquad
  C(\alpha):=\frac{4}{\left(\int_0^\alpha\sqrt{(\Finv)'(\tau)}\,d\tau\right)^{2}}
  \label{eq:thm-scaling}
\end{equation}
が成り立つ．ここで $\asymp$ は $N$ に関して比が正の定数で上下から挟まれることを表し，その定数は $\Finv$ のみに依存する（注意\ref{rem:assumption}(ii)）．
$C(\alpha)$ は $N$ に依らない正定数であり，右辺の積分は $\beta<2$ により有限である（補題\ref{lem:zador-cond}）．
すなわち，改善比の $N$ に関する冪指数は $\beta-1$ である．
\end{theorem}

定理\ref{thm:alpha} は $\alpha$ を固定した $N$ 漸近の主張である．$\alpha$ への依存は定数 $C(\alpha)$ に集約されており，その $\alpha\downarrow0$ での挙動は次の命題で厳密に与えられる．

\begin{proposition}[定数 $C(\alpha)$ の $\alpha\downarrow0$ 漸近]
\label{prop:alpha}
定義\ref{def:regvar} の仮定のもと（$\beta\in[1,2)$），$\alpha\downarrow0$ で
\begin{equation}
  \left(\int_0^\alpha\sqrt{(\Finv)'(\tau)}\,d\tau\right)^{2}
  \;\sim\;\frac{\alpha^{\,2-\beta}\,\ell(\alpha)}{(1-\beta/2)^{2}},
  \qquad
  C(\alpha)\;\sim\;4\Bigl(1-\frac\beta2\Bigr)^{2}\,\alpha^{-(2-\beta)}\,\ell(\alpha)^{-1}.
  \label{eq:alpha-scaling}
\end{equation}
\end{proposition}

命題\ref{prop:alpha} は $C(\alpha)$ 単独の漸近であり，$N$ と $\alpha$ を同時に動かす主張ではない．
共同極限を扱うには，歪み最適配置側の Zador 漸近（補題\ref{lem:zador}）が $\alpha$ について一様に成り立つことが必要である．
これは Graf--Luschgy の定理（固定分布に対する極限定理）からは自動的には従わないため，以下では仮定として明示する．
$X_\alpha:=\Finv(\alpha U)$（$U\sim\mathrm{Unif}(0,1)$），その密度を $f_\alpha$，$N$ 点 $L^1$ 最適量子化誤差を $\eN(X_\alpha)$ と書く（定義は §\ref{sec:aux}）．

\begin{quote}
\textbf{仮定 (H)}（Zador 漸近の $\alpha$ 一様性）．ある $\alpha_0>0$ に対し
\begin{equation*}
  \sup_{\alpha\in(0,\alpha_0]}\left|\frac{4N\,\eN(X_\alpha)}{\bigl(\int f_\alpha^{1/2}\,dx\bigr)^{2}}-1\right|\;\longrightarrow\;0\qquad(N\to\infty).
\end{equation*}
\end{quote}

\begin{corollary}[共同極限（仮定 (H) のもとで）]
\label{cor:joint}
定義\ref{def:regvar} の仮定が $\beta\in(1,2)$ で成り立ち，さらに仮定 (H) が成り立つとする．
$N\to\infty$，$\alpha_N\downarrow0$，$N\alpha_N\to\infty$ とすると
\begin{equation}
  \frac{\Eunif}{\Eopt}
  \;\asymp\;N^{\beta-1}\,\ell(1/N)\;\alpha_N^{-(2-\beta)}\,\ell(\alpha_N)^{-1},
  \label{eq:conj-scaling2}
\end{equation}
ここで $\asymp$ の定数は $\Finv$ のみに依存する．すなわち，冪の成分の比は $N^{\beta-1}/\alpha_N^{2-\beta}$ のオーダーである．
\end{corollary}

\begin{example}[仮定 (H) が成り立つ場合]
\label{ex:powertail}
下側裾が厳密な冪則 $\Finv(\tau)=-c\,\tau^{-(\beta-1)}$（$\tau\in(0,\alpha_0]$，$c>0$）である場合，
$X_\alpha=\Finv(\alpha U)=\alpha^{-(\beta-1)}X_1'$（$X_1':=-cU^{-(\beta-1)}$）というスケール関係が成り立つ．
$L^1$ 量子化誤差は $\eN(aX)=a\,\eN(X)$，また $\int f_{aX}^{1/2}dx=a^{1/2}\int f_X^{1/2}dx$ を満たすから，
比 $4N\eN(X_\alpha)/(\int f_\alpha^{1/2})^2$ は $\alpha$ に依存せず，(H) は $X_1'$ 単独の Zador 漸近に帰着して成立する．
一般の緩変動 $\ell$ に対する (H) の成否は本稿では扱わず，今後の課題とする．
\end{example}

\begin{remark}[$\beta=1$（ガウス裾）の扱い]
\label{rem:beta1}
$\beta=1$ では，補題\ref{lem:tail} の $\ell_1(x)=\int_x^\alpha s^{-1}\ell(s)\,ds$ を用いると，固定 $\alpha$ で
\begin{equation*}
  \Eunif\;\asymp\;\frac1\alpha\,\frac{\ell_1(1/N)}{N},\qquad
  \Eopt\;\sim\;\frac1{4\alpha N}\left(\int_0^\alpha\sqrt{(\Finv)'(\tau)}\,d\tau\right)^{2},\qquad
  \frac{\Eunif}{\Eopt}\;\asymp\;\frac{4\,\ell_1(1/N)}{\bigl(\int_0^\alpha\sqrt{(\Finv)'}\bigr)^{2}}
\end{equation*}
となる（証明は補題\ref{lem:unif}・\ref{lem:opt} と同様で，補題\ref{lem:tail} の $\beta=1$ の場合を用いる）．
$\ell_1(1/N)$ は一般に $N$ とともに（対数的に）増大するため，比を通常の意味で $\Theta(1/\alpha)$ と書くことはできない．
ガウス裾では $\ell_1(x)\sim\sqrt{2\ln(1/x)}$ であり，固定 $\alpha$ では比に $\sqrt{\log N}$ の因子が残る．
$\alpha\downarrow0$ の付記として，命題\ref{prop:alpha}（$\beta=1$）より $\bigl(\int_0^\alpha\sqrt{(\Finv)'}\bigr)^2\sim4\alpha\ell(\alpha)$，
すなわち $C(\alpha)\sim1/(\alpha\ell(\alpha))$ である．
以上より $\beta=1$ は主定理とは別に，この補正を明示した形で扱う（主定理を $\beta\in(1,2)$ に限定するのが簡潔である）．
\end{remark}

\begin{remark}[証明の対象と実験の対応]
\label{rem:scope}
定理\ref{thm:alpha} が評価するのは式\eqref{eq:bound-cvar2}の\textbf{上界汎関数}（絶対誤差積分 $I_\alpha$ に基づく量）の配置間の比である．
一方，実験Aで歪み最適配置 M2 が示す $N^{-2}$ の減衰は，$\CVaR$ の\textbf{符号付き誤差}（各区間の誤差の相殺を含む）の二次精度
（中間レポート 注意4.1）に対応し，本定理の上界（絶対誤差積分）とは別の数学的対象である．
定理は実 $\CVaR$ 誤差が同じ比だけ改善することは示さず，上界汎関数の比較にとどまる．
また，実験の M2 は正則化歪み $g_{\alpha,\lambda}$（$\lambda=0.05$）を用いた実装であり，定理の対象は $\lambda=0$ の理想配置（式\eqref{eq:def-opt}）である．
実測では，歪み最適配置の上界（絶対誤差積分）の傾きは $N^{-1}$（Zador 理論と一致），符号付き誤差の傾きは $N^{-1.86}$（二次減衰と整合）であり，
両者は明確に区別される．
\end{remark}

証明の構成は次のとおりである．§\ref{sec:aux} で補助結果（Karamata の定理，分位点関数の下側裾での漸近形，中点則の $C^1$ 版，Zador の定理）を用意し，
§\ref{sec:eval} で一様配置の上界評価（補題\ref{lem:unif}）と歪み最適配置の上界評価（補題\ref{lem:opt}）を与え，
§\ref{sec:proofs} でそれらの比をとって定理\ref{thm:alpha}，命題\ref{prop:alpha}，系\ref{cor:joint} を証明する．

\section{補助補題}
\label{sec:aux}

正則変動に関する標準的事実は Bingham--Goldie--Teugels（1987, \textit{Regular Variation}，以下 BGT）を引用する．
以下では原点 $x\downarrow0$ での正則変動を扱うが，$x\mapsto1/x$ の置換により BGT の無限遠での結果がそのまま適用できる．

\begin{lemma}[Karamata の定理（正則変動関数の積分）]
\label{lem:karamata}
$\ell$ を正値・可測な緩変動関数，$a>0$ を固定した定数とする．$x\downarrow0$ で
\begin{align}
  \int_0^x s^{p}\ell(s)\,ds&\;\sim\;\frac{x^{p+1}}{p+1}\,\ell(x)&&(p>-1),\label{eq:karamata1}\\
  \int_x^{a} s^{p}\ell(s)\,ds&\;\sim\;-\frac{x^{p+1}}{p+1}\,\ell(x)&&(p<-1).\label{eq:karamata2}
\end{align}
（BGT Prop.~1.5.8 および Prop.~1.5.10．）
\end{lemma}
\begin{proof}
\eqref{eq:karamata1}を示す．$I(x)=\int_0^x s^p\ell(s)\,ds$ とおき $s=xt$ と置換すると $I(x)=x^{p+1}\int_0^1 t^p\ell(xt)\,dt$．
緩変動の定義から各 $t>0$ で $\ell(xt)/\ell(x)\to1$ であり，一様収束定理（BGT Thm~1.2.1）によりこの収束は $t\in[\delta,1]$ で一様である．
$t\in(0,\delta)$ では Potter の上界（BGT Thm~1.5.6）$\ell(xt)/\ell(x)\le C\,t^{-\epsilon}$（$0<\epsilon<p+1$，$x$ 十分小）で支配できるから，
支配収束定理より $\int_0^1 t^p\,\ell(xt)/\ell(x)\,dt\to\int_0^1 t^p\,dt=1/(p+1)$．よって $I(x)\sim x^{p+1}\ell(x)/(p+1)$．

\eqref{eq:karamata2}を示す．Potter の上界が有効な $x_0\in(0,a]$ をとり，$\int_x^a=\int_x^{x_0}+\int_{x_0}^a$ と分ける．
第2項は $x$ に依らない有限定数であり，$p<-1$ より $x^{p+1}\ell(x)\to\infty$ であるから $o(x^{p+1}\ell(x))$ である．
第1項は $s=xt$（$t\in[1,x_0/x]$）と置換し，$\ell(xt)/\ell(x)\le C\,t^{\epsilon}$（$0<\epsilon<-(p+1)$）で支配して支配収束定理を適用すると，
$\int_1^\infty t^p\,dt=-1/(p+1)$ から従う．
\end{proof}

\begin{lemma}[分位点関数の下側裾での漸近形]
\label{lem:tail}
定義\ref{def:regvar} の仮定のもと，$\beta>1$ では $\tau\downarrow0$ で
\begin{equation}
  -\Finv(\tau)\;\sim\;\frac{\tau^{-(\beta-1)}}{\beta-1}\,\ell(\tau),
  \qquad\text{特に}\quad \Finv(\tau)\to-\infty .
  \label{eq:tail}
\end{equation}
$\beta=1$ では，$\ell_1(x)=\int_x^\alpha s^{-1}\ell(s)\,ds$ とおくと，$\ell_1(x)\to\infty$ を仮定すれば $-\Finv(\tau)\sim\ell_1(\tau)$
であり，$\ell_1$ は緩変動で $\ell_1(x)/\ell(x)\to\infty$（BGT Prop.~1.5.9a）．
\end{lemma}
\begin{proof}
$\Finv(\tau)-\Finv(\alpha)=-\int_\tau^\alpha(\Finv)'(s)\,ds$ より
$-\Finv(\tau)=\int_\tau^\alpha(\Finv)'(s)\,ds-\Finv(\alpha)$．
式\eqref{eq:regvar}を代入し，$\beta>1$ では補題\ref{lem:karamata}\eqref{eq:karamata2}（$p=-\beta<-1$）を適用すると
$\int_\tau^\alpha s^{-\beta}\ell(s)\,ds\sim\tau^{-(\beta-1)}\ell(\tau)/(\beta-1)$ を得る．定数 $\Finv(\alpha)$ は発散項に吸収される．
$\beta=1$ では $\ell_1(x)\to\infty$ の仮定のもとで $\int_\tau^\alpha s^{-1}\ell(s)\,ds=\ell_1(\tau)$ が緩変動関数を定義し，
$\ell_1(x)/\ell(x)\to\infty$ となる．なお $\ell_1(x)\to\infty$ を仮定しない場合
（例えば $\ell(s)=1/\log^2(1/s)$ で $\int_0 s^{-1}\ell(s)\,ds<\infty$），分位点は有限な下端点をもち得るため，この仮定は必要である．
\end{proof}

\begin{lemma}[中点則と区間の絶対誤差（$C^1$ 版）]
\label{lem:midpoint}
$\Finv\in C^1[a,a+w]$ が区間 $[a,a+w]$ で非減少ならば，中点 $c=a+w/2$ に対し，$M=\sup_{[a,a+w]}(\Finv)'$，$m=\inf_{[a,a+w]}(\Finv)'$ として
\begin{equation}
  \frac{m}{4}\,w^2
  \;\le\;\int_a^{a+w}\bigl|\Finv(\tau)-\Finv(c)\bigr|\,d\tau
  \;\le\;\frac{M}{4}\,w^2 .
  \label{eq:midpoint-abs}
\end{equation}
\end{lemma}
\begin{proof}
$\Finv$ の非減少性から，$[a,c]$ では $\Finv(\tau)-\Finv(c)\le0$，$[c,a+w]$ では $\ge0$ である．これを用いて部分積分すると厳密な恒等式
\begin{equation*}
  \int_a^{a+w}\bigl|\Finv(\tau)-\Finv(c)\bigr|\,d\tau
  \;=\;\int_a^{c}(\Finv)'(s)\,(s-a)\,ds+\int_c^{a+w}(\Finv)'(s)\,(a+w-s)\,ds
\end{equation*}
が成り立つ．各区間で $m\le(\Finv)'\le M$ と $\int_a^c(s-a)\,ds=\int_c^{a+w}(a+w-s)\,ds=w^2/8$ を用いると，式\eqref{eq:midpoint-abs}を得る．
\end{proof}

\begin{remark}[中点則の二次精度について]
古典的な中点則 $\int_a^{a+w}\Finv-w\Finv(c)=\frac{(\Finv)''(\xi)}{24}w^3$ は $C^2$ を要するが，
本稿の主定理の仮定は $C^1$ のみであるため，これは参考情報にとどめ，証明には補題\ref{lem:midpoint} の $C^1$ 版（式\eqref{eq:midpoint-abs}）のみを用いる．
\end{remark}

\paragraph{$L^1$ 最適量子化誤差}
確率変数 $X$ と $N$ 点のコードブック $Q=\{q_1,\dots,q_N\}\subset\R$ に対し，$N$ 点 $L^1$ 最適量子化誤差を
\begin{equation}
  \eN(X):=\inf_{|Q|\le N}\E\min_{q\in Q}|X-q|
  \label{eq:def-eN}
\end{equation}
と定義する．$\E|X|<\infty$ ならば下限は達成される（Graf--Luschgy 2000, Thm~4.12）．

\begin{lemma}[1次元 $L^1$ 最適量子化の漸近（Zador）]
\label{lem:zador}
確率変数 $X$ が絶対連続な分布をもち，密度を $f$ とする．ある $\delta>0$ に対し $\E|X|^{1+\delta}<\infty$ とする．このとき
\begin{equation}
  \eN(X)\;\sim\;\frac1{4N}\left(\int f(x)^{1/2}\,dx\right)^{2}
  \qquad(N\to\infty)
  \label{eq:zador}
\end{equation}
が成り立つ．ただし右辺の積分が発散する場合は $N\eN(X)\to\infty$ と読む．
\end{lemma}
\begin{proof}
Graf--Luschgy（2000, \textit{Foundations of Quantization for Probability Distributions}, LNM~1730）の Theorem~6.2 を $d=1$，$r=1$ で適用する．
同定理は $\lim_N N\,\eN(X)=Q_1([0,1])\,\|f\|_{1/2}$ を与え，$1$ 次元 $L^1$ の量子化係数は $Q_1([0,1])=\tfrac14$，
$\|f\|_{1/2}=\bigl(\int f^{1/2}\bigr)^2$ である．
\end{proof}

\begin{remark}[点密度の直観（ヒューリスティック）]
\label{rem:zador-heuristic}
式\eqref{eq:zador}の最適点密度は $\lambda(x)\propto f(x)^{1/2}$ で与えられる．これを分位点空間 $\tau=F(x)$（$d\tau=f(x)\,dx$）へ変換すると
$p(\tau)=\lambda(x)\frac{dx}{d\tau}=\frac{f^{1/2}}{f}=f^{-1/2}(\Finv(\tau))=\bigl((\Finv)'(\tau)\bigr)^{1/2}$ となる．
この変分計算による直観は，厳密な証明（補題\ref{lem:zador} の引用）とは区別し，点密度の解釈として記すにとどめる．
数値検証（自由度 $3$ の $t$分布での実測と Zador 式の一致，実測/理論 $\approx1.06$）は証明の一部ではなく，理論との整合性確認である．
\end{remark}

\section{一様配置の上界評価と歪み最適配置の上界評価}
\label{sec:eval}

\begin{lemma}[一様配置による上界の評価]
\label{lem:unif}
定義\ref{def:regvar} の仮定のもと，$\beta\in(1,2)$ で
\begin{equation}
  \Eunif\;\asymp\;\frac1\alpha\,N^{-(2-\beta)}\,\ell(1/N)
  \qquad(N\to\infty).
  \label{eq:eunif}
\end{equation}
さらに，$\asymp$ の定数は次の意味で $\alpha$ に依存しない．$\Finv$ のみに依存する $c_1,c_2>0$ と $h_0>0$ が存在し，
$h=1/N\le h_0$ かつ $|\Finv(\alpha)|\le\frac12|\Finv(2h)|$ を満たすすべての $(\alpha,N)$ に対して
$c_1\,h^{2-\beta}\ell(h)\le\Iunif\le c_2\,h^{2-\beta}\ell(h)+h^2(\Finv)'(\alpha/2)$ が成り立つ．
\end{lemma}
\begin{proof}
$\Iunif=\int_0^\alpha|\Finv-\Fhinv|\,d\tau$ を区間ごとに分解する．$h=1/N$，$M=\lfloor\alpha N\rfloor$ とおく．
区間 $J_i=[ih,(i+1)h]$ が $[0,\alpha]$ に含まれるのは $i=0,\dots,M-1$ であり，$\alpha N$ が整数でなければ部分区間 $[Mh,\alpha]$（幅 $<h$）が残る．
$\Finv$ は非減少で $\tau\downarrow0$ で $-\infty$ に発散する（補題\ref{lem:tail}）．

\noindent\textbf{第1区間 $J_0=[0,h]$}
原子は $q_0=\Finv(h/2)$ であり，誤差は特異性側（$\tau<h/2$）に支配される．
\begin{itemize}
\item \textbf{下界}：$[0,h/4]$ 上で $\Finv(h/2)-\Finv(\tau)\ge\Finv(h/2)-\Finv(h/4)$ であるから
\begin{equation*}
\int_0^{h}\bigl|\Finv(\tau)-\Finv(h/2)\bigr|\,d\tau
\;\ge\;\frac{h}{4}\int_{h/4}^{h/2}(\Finv)'(s)\,ds .
\end{equation*}
式\eqref{eq:regvar}より $h$ 十分小で $[h/4,h/2]$ 上 $(\Finv)'(s)\ge(1-\varepsilon)s^{-\beta}\ell(s)$ であり，
一様収束定理を固定比区間に適用して（$s=hu$）
$\int_{h/4}^{h/2}s^{-\beta}\ell(s)\,ds\sim h^{1-\beta}\ell(h)\int_{1/4}^{1/2}u^{-\beta}\,du$．
よって第1区間の寄与 $I_0$ は $I_0\ge c\,h^{2-\beta}\ell(h)$（$c$ は $\Finv$ のみに依存）．
\item \textbf{上界}：$I_0\le\int_0^h|\Finv(\tau)|\,d\tau+h\,|\Finv(h/2)|$．
補題\ref{lem:tail}と補題\ref{lem:karamata}\eqref{eq:karamata1}（$p=-(\beta-1)>-1$）より第1項は $\sim\frac{h^{2-\beta}\ell(h)}{(\beta-1)(2-\beta)}$，
第2項は $\sim\frac{(h/2)^{1-\beta}h\,\ell(h)}{\beta-1}$ であり，いずれも $\asymp h^{2-\beta}\ell(h)$．
\end{itemize}
したがって $I_0\asymp h^{2-\beta}\ell(h)=N^{-(2-\beta)}\ell(1/N)$ であり，定数は $\alpha$ に依存しない．

\noindent\textbf{正則部（区間 $i=1,\dots,M-1$）}
各 $J_i$ は $(0,\alpha]$ に含まれ原子 $\Finv((i+\frac12)h)$ も $J_i$ 上の値であるから，補題\ref{lem:midpoint}と $(\Finv)'$ の非増加性より，区間 $i$ の寄与 $e_i$ は
\begin{equation*}
  \frac{h^2}{4}(\Finv)'(\tau_{i+1})\;\le\;e_i\;\le\;\frac{h^2}{4}(\Finv)'(\tau_i)
\end{equation*}
を満たす．非増加関数 $\varphi=(\Finv)'$ に対し $h\varphi(\tau_{i+1})\ge\int_{\tau_{i+1}}^{\tau_{i+2}}\varphi$，$h\varphi(\tau_i)\le\int_{\tau_{i-1}}^{\tau_i}\varphi$（$i\ge2$）であるから，
\begin{align*}
\sum_{i=1}^{M-1}e_i
&\;\ge\;\frac h4\int_{2h}^{(M+1)h}\varphi\;\ge\;\frac h4\bigl(\Finv(\alpha)-\Finv(2h)\bigr),\\
\sum_{i=1}^{M-1}e_i
&\;\le\;\frac{h^2}{4}\varphi(h)+\frac h4\int_{h}^{(M-1)h}\varphi\;\le\;\frac{h^2}{4}\varphi(h)+\frac h4\bigl(\Finv(\alpha)-\Finv(h)\bigr).
\end{align*}
補題\ref{lem:tail}より $-\Finv(2h)\sim\frac{(2h)^{1-\beta}\ell(h)}{\beta-1}$，$-\Finv(h)\sim\frac{h^{1-\beta}\ell(h)}{\beta-1}$，
また $h^2\varphi(h)\sim h^{2-\beta}\ell(h)$ であり，$\Finv(\alpha)$ は $h\downarrow0$ で低次項である．
特に $|\Finv(\alpha)|\le\frac12|\Finv(2h)|$ のとき下界は $\frac h8|\Finv(2h)|$ 以上となり，正則部の和は $\asymp h^{2-\beta}\ell(h)$ で，定数は $\alpha$ に依存しない．
$\beta-2=-(2-\beta)$ より，正則部の和は第1区間 $I_0$ と\textbf{同じオーダー} $N^{-(2-\beta)}\ell(1/N)$ である．

\noindent\textbf{部分区間 $[Mh,\alpha]$}
$\alpha N$ が整数でない場合，幅 $w=\alpha-Mh<h$ の部分区間の原子は $q=\Finv(Mh+h/2)$ であり，これは部分区間の外に出得る．
$\tau\in[Mh,\alpha]$ に対し $|\Finv(\tau)-q|\le h\,\sup_{[Mh,Mh+h]}(\Finv)'\le h\,(\Finv)'(Mh)$ であるから，寄与は
$\le w\,h\,(\Finv)'(Mh)\le h^2(\Finv)'(\alpha/2)$（$h\le\alpha/2$ のとき $Mh\ge\alpha-h\ge\alpha/2$ と非増加性による）．
これは $O_\alpha(h^2)=o(h^{2-\beta}\ell(h))$ であり主要オーダーに影響しない．

以上より $\Iunif\asymp N^{-(2-\beta)}\ell(1/N)$ であり，式\eqref{eq:bound-cvar2}の $\frac1\alpha$ を掛けて式\eqref{eq:eunif}を得る．
\end{proof}

\begin{lemma}[歪み最適配置と $L^1$ 最適量子化の同一視]
\label{lem:identify}
$\Finv\in L^1[0,\alpha]$ とし，$X_\alpha=\Finv(\alpha U)$（$U\sim\mathrm{Unif}(0,1)$）とおく．このとき
\begin{equation}
  \Iopt\;=\;\min_{\hat Z_N}I_\alpha(\hat Z_N)\;=\;\alpha\,\eN(X_\alpha).
  \label{eq:identify}
\end{equation}
\end{lemma}
\begin{proof}
\textbf{下界}：任意の $N$ 原子分布 $\hat Z_N$ の原子を $Q=\{q_1,\dots,q_N\}$ とすると，各 $\tau$ で $\Fhinv(\tau)\in Q$ であるから
$|\Finv(\tau)-\Fhinv(\tau)|\ge\min_{q\in Q}|\Finv(\tau)-q|$．両辺を $[0,\alpha]$ で積分し $\tau=\alpha u$ と置換すると
$I_\alpha(\hat Z_N)\ge\int_0^\alpha\min_{q}|\Finv(\tau)-q|\,d\tau=\alpha\,\E\min_q|X_\alpha-q|\ge\alpha\,\eN(X_\alpha)$．

\textbf{達成}：$\E|X_\alpha|=\frac1\alpha\int_0^\alpha|\Finv|<\infty$ より最適コードブック $q_1<\dots<q_N$ が存在する．
$\Finv$ は非減少であるから，最近傍写像 $\tau\mapsto\arg\min_i|\Finv(\tau)-q_i|$（同点は小さい添字）は $\tau$ について非減少であり，
各 $q_i$ に割り当てられる $\tau$ の集合は区間 $[\tau_{i-1},\tau_i)$（$0=\tau_0\le\tau_1\le\dots\le\tau_N=\alpha$）である．
$\hat Z_N^\ast$ を，$[\tau_{i-1},\tau_i)$ 上で $\Fhinv=q_i$（$i<N$），$[\tau_{N-1},1]$ 上で $\Fhinv=q_N$ となる分布
（原子 $q_i$ に質量 $\tau_i-\tau_{i-1}$，最後の原子に質量 $1-\tau_{N-1}$．質量 $0$ の原子は除いてよい）と定めると，
$I_\alpha(\hat Z_N^\ast)=\int_0^\alpha\min_q|\Finv(\tau)-q|\,d\tau=\alpha\,\eN(X_\alpha)$ で下界が達成される．
\end{proof}

\begin{lemma}[Zador の条件は定義\ref{def:regvar} から従う]
\label{lem:zador-cond}
定義\ref{def:regvar} の仮定のもと（$\beta\in[1,2)$），$X_\alpha=\Finv(\alpha U)$ は次を満たす．
(i) $X_\alpha$ は絶対連続で，密度は $f_\alpha(x)=\frac1\alpha f(x)$（$x<\Finv(\alpha)$），$f_\alpha(x)=0$（$x\ge\Finv(\alpha)$）．ここで $f=1/(\Finv)'\circ F$ は $Z$ の下側裾の密度である．
(ii) $0<\delta<(2-\beta)/(\beta-1)$（$\beta=1$ なら任意の $\delta>0$）に対し $\E|X_\alpha|^{1+\delta}<\infty$．
(iii) $\displaystyle\int f_\alpha(x)^{1/2}\,dx=\frac1{\sqrt\alpha}\int_0^\alpha\sqrt{(\Finv)'(\tau)}\,d\tau<\infty$．
\end{lemma}
\begin{proof}
(i) $(\Finv)'>0$ より $\Finv$ は $(0,\alpha]$ から $(-\infty,\Finv(\alpha)]$ への $C^1$ 級の狭義単調増大な全単射であり，
$\Prob(X_\alpha\le x)=\Prob(\alpha U\le F(x))=F(x)/\alpha$（$x<\Finv(\alpha)$）．微分して $f_\alpha=f/\alpha$，$f(x)=1/(\Finv)'(F(x))$．
(ii) $\E|X_\alpha|^{1+\delta}=\frac1\alpha\int_0^\alpha|\Finv(\tau)|^{1+\delta}\,d\tau$．$\beta>1$ では補題\ref{lem:tail}より被積分関数は
$\tau^{-(\beta-1)(1+\delta)}\ell(\tau)^{1+\delta}$ に漸近し，$(\beta-1)(1+\delta)<1$ のとき補題\ref{lem:karamata}\eqref{eq:karamata1}より原点で可積分．
$\beta<2$ よりこのような $\delta>0$ が存在する．$\beta=1$ では $|\Finv|\sim\ell_1$ が緩変動なので任意の $\delta$ で可積分．
(iii) $x=\Finv(\tau)$（$dx=(\Finv)'(\tau)\,d\tau$）と置換すると
$\int f_\alpha^{1/2}dx=\frac1{\sqrt\alpha}\int_0^\alpha f(\Finv(\tau))^{1/2}(\Finv)'(\tau)\,d\tau=\frac1{\sqrt\alpha}\int_0^\alpha\sqrt{(\Finv)'(\tau)}\,d\tau$．
被積分関数は $\tau^{-\beta/2}\ell(\tau)^{1/2}$ に漸近し，$\beta/2<1$ より補題\ref{lem:karamata}\eqref{eq:karamata1}で可積分．
\end{proof}

\begin{lemma}[歪み最適配置による上界の評価]
\label{lem:opt}
定義\ref{def:regvar} の仮定のもと（$\beta\in[1,2)$），固定 $\alpha$ に対し
\begin{equation}
  \Eopt\;\sim\;\frac1{4\alpha N}\left(\int_0^\alpha\sqrt{(\Finv)'(\tau)}\,d\tau\right)^{2}
  \qquad(N\to\infty).
  \label{eq:eopt}
\end{equation}
\end{lemma}
\begin{proof}
補題\ref{lem:identify}より $\Eopt=\frac1\alpha\Iopt=\eN(X_\alpha)$．補題\ref{lem:zador-cond}より $X_\alpha$ は補題\ref{lem:zador}の条件を満たし，
\begin{equation*}
\eN(X_\alpha)\;\sim\;\frac1{4N}\left(\int f_\alpha^{1/2}\,dx\right)^{2}
=\frac1{4N}\cdot\frac1\alpha\left(\int_0^\alpha\sqrt{(\Finv)'(\tau)}\,d\tau\right)^{2}
\end{equation*}
を得る．これが式\eqref{eq:eopt}である．
なお，自由度 $3$ の $t$分布での数値照合（$N=256$ で理論値 $4.33\times10^{-4}$，実測 $4.60\times10^{-4}$，比 $1.06$）は
理論との整合性確認であり，証明の一部ではない（注意\ref{rem:scope}）．
\end{proof}

\section{定理\ref{thm:alpha}，命題\ref{prop:alpha}，系\ref{cor:joint} の証明}
\label{sec:proofs}

\begin{proof}[定理\ref{thm:alpha} の証明]
補題\ref{lem:unif}（式\eqref{eq:eunif}）と補題\ref{lem:opt}（式\eqref{eq:eopt}）の比をとると，固定 $\alpha$ で $N\to\infty$ により
\begin{equation*}
  \frac{\Eunif}{\Eopt}
  \;\asymp\;\frac{\frac1\alpha N^{-(2-\beta)}\ell(1/N)}{\frac1{4\alpha N}\left(\int_0^\alpha\sqrt{(\Finv)'}\right)^{2}}
  \;=\;\frac{4\,N^{\beta-1}\ell(1/N)}{\left(\int_0^\alpha\sqrt{(\Finv)'}\right)^{2}}
  \;=\;C(\alpha)\,N^{\beta-1}\ell(1/N)
\end{equation*}
を得る．$\asymp$ の定数は補題\ref{lem:unif}のものであり $\Finv$ のみに依存する．これが式\eqref{eq:thm-scaling}である．
\end{proof}

\begin{proof}[命題\ref{prop:alpha} の証明]
$\ell^{1/2}$ も緩変動である．式\eqref{eq:regvar}より $\sqrt{(\Finv)'(\tau)}\sim\tau^{-\beta/2}\ell(\tau)^{1/2}$ であり，
補題\ref{lem:karamata}\eqref{eq:karamata1}（$p=-\beta/2>-1$）を $\alpha\downarrow0$ で適用すると
\begin{equation*}
\int_0^\alpha\sqrt{(\Finv)'(\tau)}\,d\tau
\;\sim\;\int_0^\alpha\tau^{-\beta/2}\ell(\tau)^{1/2}\,d\tau
\;\sim\;\frac{\alpha^{1-\beta/2}}{1-\beta/2}\,\ell(\alpha)^{1/2}.
\end{equation*}
これを2乗して式\eqref{eq:alpha-scaling}の第1式を，さらに $C(\alpha)=4/(\int_0^\alpha\sqrt{(\Finv)'})^2$ に代入して第2式を得る．
\end{proof}

\begin{proof}[系\ref{cor:joint} の証明]
\textbf{歪み最適配置側}．補題\ref{lem:identify}より $\Eopt=\eN(X_{\alpha_N})$．仮定 (H) と補題\ref{lem:zador-cond}(iii) より，$N\to\infty$ で
$\Eopt=\bigl(1+o(1)\bigr)\frac1{4\alpha_N N}\bigl(\int_0^{\alpha_N}\sqrt{(\Finv)'}\bigr)^2$ であり，
$\alpha_N\downarrow0$ より命題\ref{prop:alpha}を適用して
$\Eopt=\bigl(1+o(1)\bigr)\frac{\alpha_N^{1-\beta}\ell(\alpha_N)}{4(1-\beta/2)^2N}$ を得る．

\textbf{一様配置側}．$N\alpha_N\to\infty$ より $\alpha_N/(2/N)\to\infty$ であり，$|\Finv|$ は原点で指数 $-(\beta-1)<0$ の正則変動関数（補題\ref{lem:tail}）であるから，
Potter の上界により $|\Finv(\alpha_N)|/|\Finv(2/N)|\to0$．よって $N$ 十分大で補題\ref{lem:unif}の条件 $|\Finv(\alpha_N)|\le\frac12|\Finv(2/N)|$ が満たされ，
$\Iunif\ge c_1N^{-(2-\beta)}\ell(1/N)$．上界側の付加項 $h^2(\Finv)'(\alpha_N/2)$ は，$(\Finv)'(\alpha_N/2)\asymp\alpha_N^{-\beta}\ell(\alpha_N)$ と
$N\alpha_N\to\infty$ から $h^2\alpha_N^{-\beta}\ell(\alpha_N)=h^{2-\beta}\ell(h)\cdot(h/\alpha_N)^{\beta}\,\ell(\alpha_N)/\ell(h)=o(h^{2-\beta}\ell(h))$
（Potter の上界による）であり吸収される．したがって $\Eunif\asymp\frac1{\alpha_N}N^{-(2-\beta)}\ell(1/N)$ で，定数は $\Finv$ のみに依存する．

両者の比をとると
\begin{equation*}
  \frac{\Eunif}{\Eopt}
  \;\asymp\;\frac{\frac1{\alpha_N}N^{-(2-\beta)}\ell(1/N)}{\frac{\alpha_N^{1-\beta}\ell(\alpha_N)}{N}}
  \;=\;N^{\beta-1}\ell(1/N)\,\alpha_N^{-(2-\beta)}\ell(\alpha_N)^{-1}
\end{equation*}
となり，式\eqref{eq:conj-scaling2}を得る．
\end{proof}

\end{document}
